\begin{ex}
Kiểm đếm số xe đi qua cổng trường trong 10 phút: 12 xe đạp, 25 xe máy, 8 xe ô tô. Có bao nhiêu xe máy đi qua?
\shortans{25}
\loigiai{
Có 25 xe máy.
}
\end{ex}

\begin{ex}
Tổng số xe các loại đi qua cổng trường trong 10 phút trên là bao nhiêu xe?
\shortans{45}
\loigiai{
$12 + 25 + 8 = 45$ xe.
}
\end{ex}

\begin{ex}
Loại phương tiện nào đi qua ít nhất? (Xe đạp ghi 1, xe máy ghi 2, ô tô ghi 3).
\shortans{3}
\loigiai{
Ô tô có 8 xe (ít nhất).
}
\end{ex}

\begin{ex}
Xe máy nhiều hơn xe đạp bao nhiêu chiếc?
\shortans{13}
\loigiai{
$25 - 12 = 13$ xe.
}
\end{ex}

\begin{ex}
Một nhóm học sinh đo chiều cao của 4 bạn: An: 132 cm, Bình: 135 cm, Cường: 130 cm, Dũng: 138 cm. Bạn cao nhất cao bao nhiêu xăng-ti-mét?
\shortans{138}
\loigiai{
Dũng cao nhất: 138 cm.
}
\end{ex}

\begin{ex}
Bạn thấp nhất cao bao nhiêu xăng-ti-mét?
\shortans{130}
\loigiai{
Cường thấp nhất: 130 cm.
}
\end{ex}

\begin{ex}
Bạn cao nhất cao hơn bạn thấp nhất bao nhiêu xăng-ti-mét?
\shortans{8}
\loigiai{
$138 - 130 = 8$ cm.
}
\end{ex}

\begin{ex}
Trong 4 bạn trên, có bao nhiêu bạn có chiều cao từ 132 cm trở lên?
\shortans{3}
\loigiai{
An (132), Bình (135), Dũng (138): có 3 bạn.
}
\end{ex}

\begin{ex}
Kiểm đếm số lượng hoa quả mẹ mua: 10 quả cam, 15 quả quýt, 5 quả xoài. Quýt nhiều hơn xoài bao nhiêu quả?
\shortans{10}
\loigiai{
$15 - 5 = 10$ quả.
}
\end{ex}

\begin{ex}
Tổng số quả mẹ mua là bao nhiêu quả?
\shortans{30}
\loigiai{
$10 + 15 + 5 = 30$ quả.
}
\end{ex}

\begin{ex}
Tung đồng xu 20 lần, thấy có 12 lần xuất hiện mặt Ngửa. Hỏi có bao nhiêu lần xuất hiện mặt Sấp?
\shortans{8}
\loigiai{
$20 - 12 = 8$ lần.
}
\end{ex}

\begin{ex}
Gieo xúc xắc 30 lần, thấy có 6 lần xuất hiện mặt 6 chấm. Số lần xuất hiện các mặt còn lại là bao nhiêu lần?
\shortans{24}
\loigiai{
$30 - 6 = 24$ lần.
}
\end{ex}

\begin{ex}
Bảng điểm kiểm tra môn Toán của tổ 1: có 3 điểm 10, 4 điểm 9, 2 điểm 8, 1 điểm 7. Tổ 1 có bao nhiêu học sinh?
\shortans{10}
\loigiai{
$3 + 4 + 2 + 1 = 10$ học sinh.
}
\end{ex}

\begin{ex}
Có bao nhiêu bạn đạt điểm 9 trở lên trong tổ 1?
\shortans{7}
\loigiai{
$3 + 4 = 7$ bạn.
}
\end{ex}

\begin{ex}
Điểm số có nhiều bạn đạt được nhất là điểm mấy?
\shortans{9}
\loigiai{
Điểm 9 có 4 bạn đạt (nhiều nhất).
}
\end{ex}

\begin{ex}
Điểm số có ít bạn đạt được nhất là điểm mấy?
\shortans{7}
\loigiai{
Điểm 7 chỉ có 1 bạn đạt.
}
\end{ex}

\begin{ex}
Theo dõi số ngày mưa trong tháng 7 ở một địa phương: có 12 ngày mưa. Hỏi có bao nhiêu ngày không mưa trong tháng 7?
\shortans{19}
\loigiai{
Tháng 7 có 31 ngày: $31 - 12 = 19$ ngày không mưa.
}
\end{ex}

\begin{ex}
Số lượng sách các môn trong tủ sách lớp: Toán: 45 cuốn, Tiếng Việt: 50 cuốn, Tự nhiên và Xã hội: 30 cuốn. Tổng số sách là bao nhiêu cuốn?
\shortans{125}
\loigiai{
$45 + 50 + 30 = 125$ cuốn.
}
\end{ex}

\begin{ex}
Sách Tiếng Việt nhiều hơn sách Tự nhiên và Xã hội bao nhiêu cuốn?
\shortans{20}
\loigiai{
$50 - 30 = 20$ cuốn.
}
\end{ex}

\begin{ex}
Trong hộp có 8 viên kẹo đỏ và 12 viên kẹo xanh. Lấy ra ngẫu nhiên 1 viên kẹo, có thể lấy được kẹo màu vàng không? (Có ghi 1, Không ghi 0).
\shortans{0}
\loigiai{
Không thể có kẹo vàng.
}
\end{ex}

\begin{ex}
Tung con xúc xắc 1 lần, khả năng xuất hiện mặt có số chấm lớn hơn 6 là mấy? (Ghi 0 nếu không thể, ghi 1 nếu chắc chắn).
\shortans{0}
\loigiai{
Không thể có mặt lớn hơn 6.
}
\end{ex}

\begin{ex}
Một bình chứa 5 quả cầu đỏ, 5 quả cầu xanh. Cần lấy ít nhất mấy quả cầu để chắc chắn có quả cầu xanh?
\shortans{6}
\loigiai{
$5 + 1 = 6$ quả.
}
\end{ex}

\begin{ex}
Đo nhiệt độ lúc 12 giờ trưa trong 3 ngày liên tiếp ghi được: $32^\circ$C, $35^\circ$C, $33^\circ$C. Nhiệt độ cao nhất trong 3 ngày là bao nhiêu độ C?
\shortans{35}
\loigiai{
$35^\circ$C là cao nhất.
}
\end{ex}

\begin{ex}
Chênh lệch giữa ngày nóng nhất và ngày mát nhất trong 3 ngày đó là bao nhiêu độ C?
\shortans{3}
\loigiai{
$35 - 32 = 3^\circ$C.
}
\end{ex}

\begin{ex}
Theo bảng số liệu rác thải nhựa thu gom được: Khối 3: 40 kg, Khối 4: 55 kg, Khối 5: 65 kg. Khối 5 thu gom nhiều hơn Khối 3 bao nhiêu ki-lô-gam?
\shortans{25}
\loigiai{
$65 - 40 = 25$ kg.
}
\end{ex}
